\section{Introducción}

\subsection{Objetivo del manual}
\pendiente{qué se logra al terminar el manual: la plataforma completa funcionando en https://defensoria-escom.ddns.net}

\subsection{Alcance}
\pendiente{qué cubre (BD, microservicios, frontends, router, HTTPS, modelo de IA) y qué no cubre (carga del histórico, que la realiza otra persona; servicios de primer contacto y subdefensoría, que pertenecen a otros equipos)}

\subsection{A quién va dirigido y conocimientos previos}
\pendiente{administrador de sistemas con nociones de Linux, SSH, contenedores y redes}

\subsection{Convenciones del documento}
\pendiente{cómo se muestran los comandos (\texttt{\#} = root, \texttt{\$} = usuario), rutas, valores que se deben sustituir como \texttt{<IP\_BD>}; las contraseñas nunca se escriben en este documento}

\subsection{Arquitectura general del despliegue}
\pendiente{diagrama TikZ de las 3 VPS: usuario $\rightarrow$ router-nginx (Frontend) $\rightarrow$ microservicios (Backend) $\rightarrow$ PostgreSQL (BD)}

\begin{table}[H]
\centering
\caption{Servidores que componen el despliegue}
\begin{tabular}{lll}
\toprule
\textbf{Rol} & \textbf{Dirección IP} & \textbf{Contenido} \\
\midrule
Base de datos & 169.58.62.99 & defensoria-db (5432), historico-db (5433) \\
Backend & 156.67.26.73 & Microservicios Spring Boot (8082--8094) \\
Frontend & 169.58.62.111 & router-nginx (80/443) y frontends Angular \\
\bottomrule
\end{tabular}
\end{table}

\subsection{Flujo de una petición}
\pendiente{ejemplo: el quejoso envía una queja $\rightarrow$ nginx /api/quejas/ $\rightarrow$ queja-service:8084 $\rightarrow$ defensoria-db; notificaciones por correo}

\subsection{Orden recomendado de instalación}
\begin{enumerate}
  \item Preparar las tres VPS (sección \ref{sec:preparacion}).
  \item Servidor de base de datos.
  \item Compilar los artefactos en el equipo de desarrollo.
  \item Servidor backend y modelo de IA.
  \item Servidor frontend, dominio y HTTPS.
  \item Datos iniciales, verificación y puesta en operación.
\end{enumerate}
